\scanpage{5}

% loc: d, upper part
d) \deleted{\unsure{Prenons donc maintenant|Passons maintenant}} Soit $E_0$ \deleted{\unsure{=}} la sous catégorie \inserted{str.} pleine de $E$ formée des $Y\in\operatorname{Ob}E$ tels que $i(Y):Y\xrightarrow{\sim}vu(Y)$. C'est donc la plus grande des sous-catégories pleines $E'$ de $E$ satisfaisant (1), et si $E'$ est une ss-catégorie pleine de $E$, elle satisfait (1) sss $E'\subset E_0$. De m\textsuperscript{e} on définit la sous-catégorie \inserted{str.} pleine $F_0$ de $F$ \inserted{formée} des \deleted{$\unsure{Q}$} $P\in\operatorname{Ob}F$ tels que $j(P):uv(P)\xrightarrow{\sim}P$, une sous catégorie pleine $F'$ de $F$ satisfait (2) sss $F'\subset F_0$. Ceci posé, on a $F_0=\overline{u(E_0)}$, $E_0=\overline{v(F_0)}$ (i.e. $F_0$\inserted{$=$} image ess. de $u|E_0$, $E_0=$ image ess. de $v|F_0$) et $u$ et $v$ induisent sur $E_0$ et $F_0$ des équivalences inverses l'une de l'autre.

% loc: e, opening
e) Pour que \deleted{l'image essentielle} l'on \deleted{ait} $v(F)\subset E_0$, il f. et s. que pour $\forall$ \deleted{$\unsure{Q}$} $P\in\operatorname{Ob}F$, on ait \deleted{$\unsure{v(P)\to vuv(P)}$}

% loc: e, isomorphism and return arrow
\diagram{\begin{tikzcd}[ampersand replacement=\&,column sep=huge]
v(P)\arrow[r,"{i(v(P))}","{\sim}" description]\&vuv(P)\arrow[l,bend left=60,dashed,"{v(j(P))\ \text{sera alors l'inverse ??}}"']
\end{tikzcd}}

% loc: e, right margin
\margin{Ou encore, il f. et s. que $v$ se factorise en $F\xrightarrow{v_1}E_1\xrightarrow{\alpha}E$, avec $\alpha$ pl. fidèle, $v_1$ \unsure{passage|pass.} à cat. des fractions, [\uline{NB} revient à $v_1$ avec un adjoint à g. $u_1=u|E_1$ mais pas néc. $\alpha$]}

% loc: e, first left diagram
\diagram{\begin{tikzcd}[ampersand replacement=\&] E\arrow[r,"u"]\arrow[d]\&F\arrow[d]\arrow[dl,"v'"']\\E_0\arrow[r,shift left=1ex,"u_0" description,"\simeq"]\&F_0\arrow[l,shift left=1ex,"v_0"'," \simeq" description]\end{tikzcd}}

% loc: e, consequence
Dans ce cas, $E_0$ est l'image \deleted{\unsure{ess.}}\inserted{essentielle} de $v$, et l'inclusion $F_0\subset\underset{\alpha}{F}$ admet un adjoint à \deleted{\unsure{gauche}}\inserted{droite}, qui est \deleted{$\unsure{une restriction à}\ F_0$} $u_0\circ\underbrace{(v\ \text{astreint à }\unsure{E|E_0})}_{v'}$ :

% loc: e, cancelled opening of calculation
\deleted{Pour que l'on \unsure{ait}}\\
$\deleted{\operatorname{Hom}_{E_0}(X_0,v'Y)\simeq\operatorname{Hom}_{E}(X_0,vu(Y))\simeq\operatorname{Hom}_{F}(uX_0,u(Y))}$

% loc: e, cancelled diagram in left margin
\diagram{\begin{tikzcd}[ampersand replacement=\&]\deleted{E}\arrow[r,"\deleted{u}"]\arrow[d]\arrow[dr,"\deleted{\unsure{u'}}" description]\&\deleted{F}\arrow[d]\arrow[dl]\\\deleted{E_0}\&\deleted{F_0}\end{tikzcd}}

% loc: e, crossed block of Hom identities
\[
\begin{aligned}
&\deleted{\operatorname{Hom}_{F_0}(u_0v'(\underset{\in\operatorname{Ob}F}{P}),\underset{\in\operatorname{Ob}F_0}{Q_0})=\operatorname{Hom}_{F}(uv(P),Q_0)
 \underset{\text{adjonction}}{\simeq}\operatorname{Hom}_{E}(vP,vQ_0)}\\
&\deleted{\simeq\operatorname{Hom}_{E}(\underbrace{vuv(P)}_{vu_0v'(P)},v(Q_0))}\\
&\deleted{\text{car }vP\simeq vuvP\longleftarrow\ \unsure{\text{redonne}}}\\
&\deleted{\simeq\operatorname{Hom}_{E_0}(v_0u_0v'P,v_0(Q_0))}\\
&\deleted{\simeq\operatorname{Hom}_{F_0}(u_0}
\end{aligned}
\]

% loc: e, proof in brackets
[En effet, on a pour $Q_0\in\operatorname{Ob}F_0$, $P\in\operatorname{Ob}F$

% loc: e, proof display
\[
\begin{aligned}
\operatorname{Hom}_{F}(Q_0,P)&\underset{\text{car }u_0v_0\simeq\operatorname{id}}{\simeq}
\operatorname{Hom}_{F}(u_0v_0Q_0,P)\simeq\operatorname{Hom}_{E}(v_0Q_0,vP)=\operatorname{Hom}_{E_0}(v_0(Q_0),v'P)\\
&\underset{u_0\ \text{pl. fid}}{\simeq}\operatorname{Hom}_{F_0}(u_0v_0Q_0,u_0v'P)
\underset{(u_0v_0\simeq\operatorname{id})}{\simeq}\operatorname{Hom}_{F_0}(Q_0,u_0v'P)\big].
\end{aligned}
\]

% loc: lower left, reverse diagram
\diagram{\begin{tikzcd}[ampersand replacement=\&]E\arrow[d]\arrow[dr,"u'"]\&F\arrow[l,"v"']\arrow[d]\\ E_0\arrow[r,shift left=.8ex,"u_0","{\simeq}" description]\&F_0\arrow[l,shift left=.8ex,"v_0"'," {\simeq}" description]\end{tikzcd}}

% loc: lower bracket
[\unsure{De m\textsuperscript{e}}, \deleted{on trouve}\inserted{est-il vrai} que $E_0\hookrightarrow E$ admet $v'u$ comme adjoint à \deleted{\unsure{gauche}}\inserted{gauche}, ou qu'il admet un adjoint à \deleted{\unsure{gauche}}\unsure{?} ?? Revient à ceci : si $\alpha v_1$ a un adjoint à g., avec $\alpha$ pl. fid et $v_1$ passage à cat. fractions, a-t-on que $\alpha$ a un adj. à g ? ]

% loc: bottom, cancelled calculation
\[
\deleted{\operatorname{Hom}_{E}(\underset{\operatorname{Ob}E}{Y},\underset{\operatorname{Ob}E_0}{X_0})\underset{v_0u_0\simeq\operatorname{id}}{\simeq}\operatorname{Hom}_{E}(Y,v_0u_0X_0)\underset{\text{adjonction}}{\simeq}\operatorname{Hom}_{\unsure{E|F}}(uY,u_0X_0)}
\]

% loc: foot, continuation
De façon duale, pour que l'on ait $u(E)\subset F_0$, il f. et s. que l'on ait \deleted{pour} $\forall X\in\operatorname{Ob}E\quad uvu(X)\xrightarrow[\sim]{j(u(X))}u(X)$, et alors $F_0$ est l'image ess. de $E$ par

% archive sheet number: 4
